% Section 1.5: Summary and Transition to Chapter 2

\section{Summary}
\label{sec:ch1_summary}

This chapter has established the theoretical and observational foundation for the dark matter searches that follow.
Multiple independent lines of evidence --- flat rotation curves in spiral galaxies, the dynamics and gravitational lensing of galaxy clusters, and the precision cosmology of the cosmic microwave background --- converge on the $\Lambda$CDM concordance model, in which roughly 85\% of the matter density is composed of cold, non-baryonic dark matter.
Big Bang Nucleosynthesis independently confirms the non-baryonic character of this component, requiring that dark matter consist of a new form of matter beyond the Standard Model.

Among the candidates proposed to fill this role, the Weakly Interacting Massive Particle stands out because of the ``WIMP miracle'': a stable particle with electroweak-scale mass and coupling, produced thermally in the early universe, naturally yields a relic abundance consistent with observation.
The thermal freeze-out calculation fixes the annihilation cross-section at $\langle \sigma v_\mathrm{rel} \rangle_\mathrm{cosmo} \approx 2.2 \times 10^{-26}$~cm$^3$/s, providing a concrete, parameter-free experimental target in the GeV--TeV mass window bracketed by the Lee--Weinberg and unitarity bounds.
The same interaction that sets the relic density can be probed from three complementary directions --- scattering in underground detectors (direct detection), production at colliders, and annihilation in astrophysical environments (indirect detection).
Among these, indirect detection via gamma-rays occupies a unique position: it directly probes the annihilation cross-section, the quantity that determines the cosmological abundance.

The $J$-factor and $D$-factor formalism developed in this chapter factorizes the predicted gamma-ray flux into a particle physics term and an astrophysical term, providing the mathematical framework that underpins every analysis in this thesis.
The current null results of WIMP searches --- with Fermi-LAT constraints from dwarf galaxies already excluding the thermal relic cross-section for masses below ${\sim}\,100$~GeV in the $b\bar{b}$ channel --- demonstrate that the era of simple signal-over-background searches is reaching its limits.
Extracting weaker signals from noise-dominated data demands the statistical and machine learning methods developed in the remainder of this thesis.
The next chapter introduces the instrument that collects the data: the Fermi Large Area Telescope and the astrophysical gamma-ray sky that constitutes both the observational setting and the primary source of systematic uncertainty for the analyses that follow.
